\subsubsection{French Brawl}

\paragraph{Outline}

French Brawl can be both a deeply theoretical template, and luck-based brawl. Have you seen Starship Troopers? On the surface you get a glossy late-90's space marine movie of bugs exploding, but if you watch it closely you can spot the critique and messages about media indoctrination and fascism and say "aha!". The movie is fully capable of wearing two hats at once and it's versatility enables it to reach broader audiences. French Brawl is the Starship Troopers/ Robocop / Fight Club of Warzone.
\\
\\
\paragraph{Distributions with predominantly small bonuses}

These boards are all about army efficiency and revolve around specific territory thresholds and FTBs. You cannot, should not, under any circumstances pick remote single picks in slow bonuses like \href{https://www.warzone.com/MultiPlayer?GameID=37197664}{the +4} in Aquitaine. Unless you are batshit crazy like me \href{https://www.warzone.com/MultiPlayer?GameID=41429904}{here vs Jake}. Solo pick +3 bonuses can work but you can still end up in situations like \href{https://www.warzone.com/MultiPlayer?GameID=39666704}{here} where they are just too slow to matter. Best way to setup a solo +3 pick is to complete it in 2 turns with all 8/8 of your deployments and jump to 5 territories in the +3 and 5 more in your other 5 spawns -> 8 Income. If you get eliminated somewhere, that's 9 territories -> 7 Income which is easily obtainable by a 10 territories and a +2 bonus or similar configurations. As such, you want coverage in a way that both denies bonuses and doesn't get eliminated, like \href{https://www.warzone.com/MultiPlayer?GameID=41381786}{here by Aika}.
\\
\\
Outside of such niche tactics, you will get the most out of picking in clusters and adjacent bonuses, due to how those armies are useful immediately. Examples of this:
\begin{itemize}
    \item \href{https://www.warzone.com/MultiPlayer?GameID=21713173}{Beep's 12th} is wrong, 48-Lozère is useful immediately (Othello's 8th).
    \item \href{https://www.warzone.com/MultiPlayer?GameID=40224644}{Aika's 3rd} forces his hands, 48-Lozère allows more ideas.
    \item \href{https://www.warzone.com/MultiPlayer?GameID=41569033}{Adjacent picks} allow me to get Income from extra territories without deployments. Plat has to commit resources in isolated picks for income. 
    \item \href{https://www.warzone.com/MultiPlayer?GameID=41393844}{vs nick} I'm playing the game 6vs4 spawns till turn 3.
    \item \href{https://www.warzone.com/MultiPlayer?GameID=38516578}{Markus vs Aika}, their 8th pick is considerably better than their 7th pick, due to being useful T1, as seen from Markus having a passive pick.
\end{itemize}

In this broad category of distributions, the entire game takes place North-Center and the ordering of picks as well as gameplay and win conditions dictate the game. Take for example this game \href{https://www.warzone.com/MultiPlayer?GameID=36761893}{by Octane} and this game \href{https://www.warzone.com/MultiPlayer?GameID=35996112}{by Platinum} and try to envision compensation for losing first pick.
\begin{enumerate}
    \item In Octane's distribution, Auvergne (Center-South +2) FTB is good but you do not want to over-commit to picking it. If you go for it as 1/2/3, that 2/3 is so high up that will lead you into losing coverage elsewhere. On top Octane's FP is the only territory controlling both +2 FTBs. The way Octane chooses to cover the south +2 FTB in case he loses FP is by knowing that if he loses it, he can first move Turn 1 counter from his 4th. Lost FP -> 2nd pick in picks, 1st Turn 1. 
    \item In Platinum's distribution, 21- Côte-d'Or serves the same purpose in controlling both FTBs. The correct approach is to 1/2 the center as Alecto does. Losing 1 means you can still counter t1 due to move order
\end{enumerate}

\textbf{Example 1}.
You can extend this idea from setting up your 1/2 properly to setting up more picks properly. For example \href{https://www.warzone.com/MultiPlayer?GameID=37941014}{here}, reverse engineer Octane's and Ralph's methodology keeping in mind that if both players mirror, FP gets 1/4/5/8/9/12 while SP receives 2/3/6/7/10/11. Ralph's set maximizes the probability of getting the +2 FTB while attempting to nullify single picks of your opponent's through a clever 6/7 adjacent pick combo that covers the +3 and +1 bonus (and expands if needed in the +2). The 6/7 strength is even greater when considering you expect mirroring 1-4 and contact on those 2 points. The glaring issue in Ralph's set is that FP has a disadvantage if he loses 6/7, as 5/8/9/ are isolated and 12 is sort of bad; you are forced to force a coin-flip elimination south or in the +2 FTB, but SP connects adjacent picks, keeps army advantage, and plays a 10terr (5vs4 income) game. Octane's set tackles this experimentally by combo-ing 12 with 4/5 for FP and covering the +2 FTB in return with 4/5 which can be lost only as second pick (assuming mirror south). The compensation of losing 4/5 and mirroring 1/2 or 1/3 is that you always get 1/2/6/7 (or 2/3/6/7). What this means in practice:
\begin{itemize}
    \item If Octane gets FP, he gets his 1/2/4/7/8 with his 4 right between 2 isolated Ralph picks. If Ralph doesn't spot how valuable the +3 combo is, Octane as FP gets 1/2/4/6/7 that's positionally winning.
    \item If Octane doesn't get FP, his compensation is 2/3/6/7 with a guaranteed elimination south T1 (given Opponent gets FTB) and 2 spawns north to muddy the waters till he finishes the +3.
\end{itemize}

The reality is that the distribution was inherently favoring SP due to how 12th is always awkward, the point of this section is mainly how you can identify that through mere calculation and experiment around it. For the actual gameplay you need to play with tempo.

\textbf{Examples of games decided by Income from Territories }.

A handful of games fit here, such as:
\begin{enumerate}
    \item \href{https://www.warzone.com/MultiPlayer?GameID=32875754}{Xeno vs Malakkan}. The game end of t2 is 6vs6 Income (9vs10 Territories), and Xeno's approach to surviving is trying to get to 10 territories. Surviving turn 3 enables forcing moves, as t4 with a 6 attack FO, he protects his bonus unless Malakkan full deploys, guesses the right border.
    \item \href{https://www.warzone.com/MultiPlayer?GameID=32875753}{tornado vs Rufus}. The game is a delicate balance of income from territories accompanied by +1 bonuses. T2 and especially T3 tornado is better accepting the brawl and trying to retain an edge.
    \item \href{https://www.warzone.com/MultiPlayer?GameID=32726827}{Xeno vs FSG}. Xeno matches income from seemingly random territories T2. Income from territories is a lot safer than Income from bonuses cause defensive kill-rates exist.
    \item \href{https://www.warzone.com/MultiPlayer?GameID=33055938}{Octane vs nick}. Octane's T2 forces 9vs9 territories game, T3 forces 12vs11 territories game. If he doesn't match nick's tempo it's most likely a loss.
    \item \href{https://www.warzone.com/MultiPlayer?GameID=37222704}{Rufus vs Aika}. Rufus' T1 went horrible and must deny 10 territories T2 to bridge the income difference.
\end{enumerate}

As a general idea in terms of gameplay, when ahead look for forcing plays, when behind take risks in any way or form. Sometimes being bold works in your favor, like how \href{https://www.warzone.com/MultiPlayer?GameID=38662026}{here} ps finds the best T3 sequence, how \href{https://www.warzone.com/MultiPlayer?GameID=38696092}{here} Markus defends a triple border like a magician, how \href{https://www.warzone.com/MultiPlayer?GameID=37731836}{Xeno} pretends he is playing vs neutrals, or how \href{https://www.warzone.com/MultiPlayer?GameID=41449221}{here} I decide to do a bit of trolling.


%Case study of \href{https://www.warzone.com/MultiPlayer?GameID=32875754}{Xeno vs malakkan} from the seasonal ladder. A board rich in +2 FTBs, +1 FTBs and adjacent picks. Notice and mimic that in these boards you are forced to have your late picks be adjacent picks in order to have useful armies T1. On picks, Xeno's compensation for each lost combination works this way: Losing 2/3/6 -> Get 1/4/5/7/8 or 1/4/5/8/9. Uphill battle but elimination south for 5vs6 income game with 9 as a counter to the north works. Losing 1/4/6 -> 2/3/5/7 and SP. If you get 8 you have safe red +3 West and threaten the lost +2 FTB through a double border anytime. Game reaches T2 and by that time it's about identifying your win condition between Income from Bonuses and Territories. 


\paragraph{+3 FTBs/combos \& +4 FTBs/combos}

\textbf{+3 FTBs/combos}
Easier to go through examples for Bretagne, West Dark Green +3, which apply for every +3 Bonus.

\begin{itemize}
    \item Solo pick explained how it requires setup above in \href{https://www.warzone.com/MultiPlayer?GameID=41381786}{Aika vs Octane}. that gets you to either 7 or 8 income, depending on the borders, end of T2.
    \item Double pick is more flexible like how in \href{https://www.warzone.com/MultiPlayer?GameID=33100117}{Buns vs Beren} he sets up 9 income (12 terr) end of T2 with 6 deployments for the +3. After the early game, Beren has from the south 3 Income from Bonuses and 4 from 9 Territories, while Buns has 3 from Bonuses and 3 from 6 Territories - power of efficiency, at the cost of an early army disadvantage.
    \item How high you prioritize it depends solely on your gameplay and distribution. In poorer distributions such as \href{https://www.warzone.com/MultiPlayer?GameID=34380051}{here between Rufus and Octane}, the west is richer in income than the east and you ought to balance picks accordingly. \href{https://www.warzone.com/MultiPlayer?GameID=32905758}{Here}, Juan overcomes the early disadvantage in armies through safe income. Generally speaking, it's better to play with the mindset that safe income is overrated, and to test limits. LF works in your favor, you can force eliminations knowing of the +3 combo. If the board is the poorest you've seen with wastelands in all the usual FTBs, the entire game is played \href{https://www.warzone.com/MultiPlayer?GameID=38391065}{West}, though you still have to be careful of over-extending.
    \item If the rest of the board is rich in income, it's better to cover the FTBs and then the slower combo +3s, something very relative to each distribution. \href{https://www.warzone.com/MultiPlayer?GameID=34192676}{Here} Rufus commits 1-8 in +2/+1 FTBs and 9-11 to cover the +3/+4 area, while Juan similarly commits 1-7 in +2/+1 FTBs and 8-10 for the +3/+4 area. \href{https://www.warzone.com/MultiPlayer?GameID=41419436}{Here} me and Octane commit 6-7/6-8 for that combo and 1-5 for +1 FTBs, \href{https://www.warzone.com/MultiPlayer?GameID=30854262}{here} Lepanto and Buns cover the +3 combo on 8-9/9-10 and +2/+1 FTBs earlier - you get the pattern. \textbf{This isn't a general rule you should follow through, picking in French Brawl is more of an ABBA puzzle with efficient use of armies.} FTBs accelerate this, while slower combos give opportunities to your opponents to eliminate you, like \href{https://www.warzone.com/MultiPlayer?GameID=33055938}{here}, \href{https://www.warzone.com/MultiPlayer?GameID=39666704}{here} and \href{https://www.warzone.com/MultiPlayer?GameID=26328881}{here}. So the answer is to calculate.
\end{itemize}

\textbf{+4 FTBs/combos}

\begin{itemize}
    \item +4s require perfect setup. Sometimes even their threat is better than the execution.
    \item For Aquitane. Example 1. \href{https://www.warzone.com/MultiPlayer?GameID=28062679}{Rufus vs Gak}, Rufus' win condition after T1 involves a) get army advantage and b) keep the game at 5vs5 income. After T2 works in his favor, T3 gak has to gamble something towards 10 territories and 6 income without dropping a territory nor the +4. Part of the cluster's strength is that Gak has to take initiatives and gambles, while Rufus can play defensively for small advantages.
    \item Example 2, Playable Aquitane and Poitou Charentes (Red +3) 
    \begin{itemize}
        \item \href{https://www.warzone.com/MultiPlayer?GameID=26158910}{Bechaa vs mod}, The +4 in Aquitaine is a big threat and Bechaa covers the entire side of the board with two picks while mod overextends.
        \item \href{https://www.warzone.com/MultiPlayer?GameID=25460540}{Rufus vs Ursus}, threat of both the +4 and +3.
        \item \href{https://www.warzone.com/MultiPlayer?GameID=30110334}{Mori vs Yelo}, the entire region is a threat so covering the +4 > over-committing for it in picks.
    \end{itemize}
    \item For Aquitaine specifically sometimes you have a triple pick where the first pick at the +1 is game-winning. Those distributions can be more rps as FP is favored. For example, 
    \begin{enumerate}
        \item \href{https://www.warzone.com/MultiPlayer?GameID=24218428}{Plat vs Muli}, Plat has a playable position purely due to winning first pick (though both sets are confusing to me).
        \item \href{https://www.warzone.com/MultiPlayer?GameID=27625748}{Bechaa vs Rento}, Bechaa does what Plat did in the previous bulletpoint, and gets a completely lost position due to losing first pick.
        \item  1st pick the +1 and triple pick into complete by T2 for 9 Income \href{https://www.warzone.com/MultiPlayer?GameID=26081143}{by Bechaa} (MGO ignored the region)
    \end{enumerate}
   \item Aquitaine triple pick expanding into the +7 \href{https://www.warzone.com/MultiPlayer?GameID=33443166}{by Octane}.
   \item Aquitaine \&/or +7 becoming useless picks \href{https://www.warzone.com/MultiPlayer?GameID=14429415}{here}, \href{https://www.warzone.com/MultiPlayer?GameID=25537788}{here}, \href{https://www.warzone.com/MultiPlayer?GameID=34470203}{here}, \href{https://www.warzone.com/MultiPlayer?GameID=36893709}{here}. Just examples of picksets going wrong with single picks.
   \item For Centre, Purple +4 in the center, \href{https://www.warzone.com/MultiPlayer?GameID=27743610}{how to do it} by Ralph. And well, you can also get it by Turn 2, like \href{https://www.warzone.com/MultiPlayer?GameID=41393844}{here}. Very rare thing to finish the bonus but please don't ignore it in picks if it's an FTB. 
   \item Languedoc Roussillon - Light Green +4 Far South is a great triple pick, weird double pick, and awesome quadruple pick with a +7. example of it being contested by both players \href{https://www.warzone.com/MultiPlayer?GameID=32881674}{by Octane and We ride!}.
   \item Triple pick \href{https://www.warzone.com/MultiPlayer?GameID=28112223}{by Roi on 3/4/5}, \href{https://www.warzone.com/MultiPlayer?GameID=27829090}{by Guest on 4/5/6}, \href{https://www.warzone.com/MultiPlayer?GameID=24276469}{by DrApe in 4-7}
   \item Double pick \href{https://www.warzone.com/MultiPlayer?GameID=22632740}{by Rufus on 5/6}, \href{https://www.warzone.com/MultiPlayer?GameID=14530998}{by Rufus on 4/5} and then a single pick \href{https://www.warzone.com/MultiPlayer?GameID=14514975}{by mod} somehow working.
   \item Isolated pick -> Defeat -> be careful \href{https://www.warzone.com/MultiPlayer?GameID=31721614}{how you pick}.
\end{itemize}

Don't get baited by the games I've linked, the +4 bonuses have 48\%(Aquitaine) and 40\%(Light Green one) WR in games above 1800 Rating. They require specific setups and some bailing from your opponent misplaying. But, if you complete one of them and it's uncontested, 4 Income for 6 Territories is a cheat-code. In comparison, 6 Territories from 2x+1 Bonuses is 2 Income. 
\\
\\
Overall, I went through basic combos and Income transitions such as:
\begin{itemize}
    \item Solo pick +3, 8/8 deployments, 4->4->7/8 (8 if 10 territories)
    \item Combo pick +3, 6/9 deployments, 4->5->9 with 10 territories
    \item Triple pick +4, 6/8 deployments, 4->4->9 with 10 territories
\end{itemize}

If you have a board with such combos, the question becomes whether you can keep up to the FTBs. A +2 FTB will go 4->7 end of t1 and immediately look to expand/eliminate/counter your bonus. So the discussion revolves around optimizing Turns 1-3.

\paragraph{+7s}

If the +4s where efficiency bombs, then wait till you hear this comparison:
\begin{itemize}
    \item 9 Territories for a +7 Bonus -> +7 Income
    \item 9 Territories for 3x +1 Bonuses -> +3 Income
\end{itemize}

\textbf{Rhone Alpes, Dark Green +7}

Side note, somehow across 400 games >1800 rating, this bonus has 65\% WR, but at lest  90\% of those games it's picked in combination to a +2 bonus nearby.
\\
\\
This bonus requires an unbelievably rare setup since you are likely looking at completing it by Turn 3 and that is enough time for someone to close the distance. Counter example, \href{https://www.warzone.com/MultiPlayer?GameID=27915506}{in this distribution} the +7 looks nice can 5x pick it, but above it there's 3 +2 FTBs. Pardon's 1/2 double borders the +7 on turn 2 -> you could only complete it on turn 3 -> gg.
\\
\\
Therefore, need 3 step distance or higher.
\\
\\
Example 2, \href{https://www.warzone.com/MultiPlayer?GameID=15116307}{here}, Limberson commits to the 6x pick but handicaps himself on Turn 3 staying at 15 territories and not using +3 to move north. At 15vs16 with 70-Haute-Saône it's a different game. And Turn 4 you cannot be passive in a 3 territory +1 bonus that's punishable by the International Court of Justice.
\\
\\
Example 3, \href{https://www.warzone.com/MultiPlayer?GameID=40231108}{here}, I commit to a +7 that's 2-3 steps away from the opponent (and misplay it) -> gg.
\\
\\
Example 4, \href{https://www.warzone.com/MultiPlayer?GameID=38011285}{by Xeno}, 4->5->11->20 with a triple pick +4 end of T2 and the +7 T3, impossible to lose.
\\
\\
\textbf{PACA, Red +7}

This one requires the stars to align. With 2 picks it's too slow unless you have an FTB, while with 3 picks by definition you need the +4 and north +7 spawn.
\\
\\
Example 1, \href{https://www.warzone.com/MultiPlayer?GameID=34479630}{by Rufus vs Xeno}, Rufus completes the +7 on T4, 4->5->9->13->22->33 which is lowkey crazy. At one point Xeno has 21 income from 42 territories and 0 bonuses.
\\
\\
Example 2, \href{https://www.warzone.com/MultiPlayer?GameID=40231108}{by 7ate9 vs Carpenter}, I should've went for the south +7 like \href{https://www.warzone.com/MultiPlayer?GameID=40274907}{this} but had no idea it was possible after picks.
\\
\\
Example 3, \href{https://www.warzone.com/MultiPlayer?GameID=41470794}{by 7ate9 vs Lumumba}, I knew this is likely troll but decided to go for it regardless. Worked out. The thing about that +7 precisely is that by the time you are broken, you can build an army advantage elsewhere, pivot, and still get the income from the territories there.
\\
\\
\textbf{Midi Pyrenees, Blue +7}

This one is the most common and powerful due to how fast it can expand to both +4s next to it.
\\
\\
Example 1. \href{https://www.warzone.com/MultiPlayer?GameID=41470788}{Xeno vs Octane in OP} and \href{https://www.warzone.com/MultiPlayer?GameID=32934772}{Xeno vs Octane in the seasonal}, of 6x commit to south and 4x commit south + 2x cover north.
\\
\\
Example 2, on pivoting to a +7 in the mid-game. \href{https://www.warzone.com/MultiPlayer?GameID=32666150}{Xeno vs mod, seasonal}, features a triple pick +7 to pivot in the middle game, \href{https://www.warzone.com/MultiPlayer?GameID=34614356}{Rufus vs Sepp}, double pick +7 for +2 FTB into pivot.
\\
\\
Example 3. \href{https://www.warzone.com/MultiPlayer?GameID=38140305}{Rufus vs Guest}, 4x +7 and cover north.
\\
\\
\textbf{How not to play with a +7 Bonus}

Step 1. Don't use Wild Strawberry as your color. Step 2, do the opposite Wild Strawberry players did in the following games, at first in \href{https://www.warzone.com/MultiPlayer?GameID=26670234}{Ralph vs The Infinite Rick} and second \href{https://www.warzone.com/MultiPlayer?GameID=41495186}{MarkusBM vs Widzsiz}. In both positions, especially on Turns 6-7 (and afterwards), you should be greedier and convert your army advantage through breaking opponent bonuses. Overcommit south trying to defend = losing condition.








%French Brawl early game is all about army efficiency. It's why combining 2 starting spawns is so good.

%\href{https://www.warzone.com/MultiPlayer?GameID=38516577}{ml3 game}
%-> beren's picks are 123 center, 456 green +1, then stack cover blue +3. 
%-> ideally would rather 1/2 center, then cover +3 cluster, then green +1 cluster
%+ adjacent picks are better than isolated picks, like 9

%\href{https://www.warzone.com/MultiPlayer?GameID=38516578}{ml3 game}
%8 better than 7 cause adjacent, armies immediately useful

%\href{https://www.warzone.com/MultiPlayer?GameID=32666145}{ruf vs jacob} on how to structure picks for combos.

%\href{https://www.warzone.com/MultiPlayer?GameID=32666150}{seasonal} on 2x +7 combo 

%\href{https://www.warzone.com/MultiPlayer?GameID=32875753}{taking neutrals} correct play on t2 is to find a way to change the position from 8/8 territories to 9/7 in your favor.

%"I see people misuse delays though a lot in FB.  Most notably by tapping unguarded opponent territories."

%\href{https://www.warzone.com/MultiPlayer?GameID=34479630}{big bonuses} the extent of income per territories 

%\href{https://www.warzone.com/MultiPlayer?GameID=34236815}{t2 misplay}. Light Green is borderline at forced win territory by end of turn 1 if they capture Picardie with +2 and first order, +1 south stay and +1 east stay.

%In very brawly FB distributions where you have to commit 12 picks north or 10/9 north and 2/3 center, it becomes a game of: 
%\begin{itemize}
    %\item guaranteeing coverage of all FTBs (especially +2s)
    %\item have useful armies immediately without deployments %through adjacent territories
    %\item try and sneak a +1 ftb in case your opponent %overdeploys in a region. similarly for stb +2.
%\end{itemize}


%\href{https://www.warzone.com/MultiPlayer?GameID=35684068}{alex's clues} on having the surrounding instead of central pick. sometimes its better to have the central pick since you can snipe either of the two outter territories with fo on turn 1 and an attack of 5/6. Also if you have center + outter, you can connect and develop, which comes at the cost of having -1 pick elsewhere. Also if opponent requires heavy deployment elsewhere, then better not to be in the center, since center requires deployments to survive.

%\href{https://www.warzone.com/MultiPlayer?GameID=35996112}{safe +3}. better way is to 1/2 center, with plat's 5 as 1, and plat's 1 as 2. losing 1 then means you can still counter t1 due to move order shenanigans. 3/4 safe combo +3. then some higher prioritising of the far north +1 picks.

%\href{https://www.warzone.com/MultiPlayer?GameID=36186163}{weirder boards}. multiple ways to tackle this one. dogpoo goes for 1-4 east, 5-9 center, 7, 10-12 north and play the brawl. awesomeusername goes for non-adjacent picks and ends up in a distribution with two practically useless starting territories. another way structure picks has mixing things up with 1 far north, 234 east, 56 north, 789 south, 10 onwards mixture of north. but all of these underestimate the combo +7 south. in a way, you want to cover that, but cannot afford to cover it very early.

%\href{https://www.warzone.com/MultiPlayer?GameID=37197664}{rule} don't single pick big bonuses (+4 or bigger). you end up with a useless spawn and lose in the brawly region due to army disadvantage.

%\href{https://www.warzone.com/MultiPlayer?GameID=38011285}{big bonus stack} 

%\href{https://www.warzone.com/MultiPlayer?GameID=40231108}{misplay analysis}. I went for a cluster that should've been played this \href{https://www.warzone.com/MultiPlayer?GameID=40274907}{way}.